% \chapter{Introduction}\label{chap:introduction}

% Since the advent of modern programming languages, developers have sought to ...

% \paragraph{Organization}  Chapter~\ref{chap:background} gives a brief overview of the background information.

% Bruzzone~\cite{Bruzzone2025papers}

% Bruzzone e Favalli

% Bruzzone et al.

% \section{Prova}

% \lipsum[1-2]

% \subsection{Prova 2}

% \lipsum[1-2]

% \paragraph{Premise} This thesis focuses on the analysis of Rust programs using static analysis techniques. Rust is a systems programming language that emphasizes safety and performance.

% \paragraph*{Problem Statement} The main problem addressed in this thesis is the challenge of ensuring memory safety in concurrent programming environments.

% \chapter{Background}\label{chap:background}




% In this chapter, we provide the necessary background to understand the
% contributions of this thesis. 


% \section{The Rust Programming Language}\label{sec:rust}

% Come vediamo nella funzione \inlinerust{fn main() {}}
% Come vediamo nella funzione \inlinerust{fn main() {}}
% \textit{Corsivo} mentre \textbf{grassetto} normale \texttt{monospaziato}.

% \mytool è uno strumento per l'analisi statica di programmi scritti in Rust.
% \slr{ciao come stai?} ciao ciao

% \lipsum[1-2]

% \begin{wrapfigure}{l}{0.5\linewidth}
% 	\centering
% 	\captionsetup{type=listing,skip=2pt}
% 	\showrust[0.45\textwidth]{ex.rs}
% 	\caption{A simple Rust program.}
% 	\label{fig:rust-hello-world}
% \end{wrapfigure}

% \lipsum[1-2]

% \begin{wrapfigure}{r}{0.5\linewidth}
% 	\centering
% 	\captionsetup{type=listing,skip=2pt}
% 	\showrust*[0.45\textwidth]{ex.rs}
% 	\caption{A simple Rust program.}
% 	\label{fig:rust-hello-world2}
% \end{wrapfigure}

% \begin{Listing}[t]
% 	\centering
% 	\showrust*[\linewidth]{ex.rs}
% 	\caption{The \texttt{SymbolTableEntry} interface.}%
% 	\label{lst:typelang:SymbolTableEntry}
% \end{Listing}

% \lipsum[1-2]
% \lipsum[1-2]
% \lipsum[1-2]



